\pagebreak
\section{Метод решения}
Ранжирование добавлено в поисковую систему ЛР7--ЛР8 так же, как стемминг (ЛР4) и сжатие (ЛР11): индекс, построенный с новым ключом, сообщает о себе флагом в заголовке, и поиск ЛР8 по этому флагу включает ранжирование. Корпус --- 1\,000\,000 документов ЛР1 (Wikipedia, Reddit, Travel Stack Exchange, Wikivoyage).

\subsection{Схема весов lnc.ltc}
Взята традиционная схема из лекции \enquote{lnc.ltc} (обозначения SMART: первая тройка букв --- веса документа, вторая --- запроса). Логарифмы десятичные.
\begin{itemize}
    \item Документ (lnc): вес терма $t$ в документе $d$ --- $1 + \log_{10} \mathit{tf}_{t,d}$, без idf, вектор нормируется на длину:
    $$w_{t,d} = \frac{1 + \log_{10} \mathit{tf}_{t,d}}{\|d\|}, \qquad \|d\| = \sqrt{\sum_{t' \in d} (1 + \log_{10} \mathit{tf}_{t',d})^2}.$$
    \item Запрос (ltc): вес терма --- $(1 + \log_{10} \mathit{tf}_{t,q}) \cdot \log_{10} \frac{N}{\mathit{df}_t}$, вектор нормируется на длину.
    \item Оценка документа --- косинус угла между векторами, т.\,е. скалярное произведение нормированных векторов:
    $$\mathrm{score}(q, d) = \sum_{t \in q \cap d} w_{t,q} \cdot w_{t,d}.$$
\end{itemize}
Почему так:
\begin{itemize}
    \item Нормировка документа обязательна для нашего корпуса: статья Wikipedia в среднем в 6 раз длиннее поста Reddit (1\,364 и 214 слов). Без нормировки длинная статья, где слово запроса встретилось много раз между делом, всегда была бы выше короткого поста, целиком посвящённого вопросу.
    \item Логарифм частоты: 10 вхождений слова не делают документ в 10 раз релевантнее одного.
    \item idf только у запроса: длину документа $\|d\|$ можно посчитать один раз при индексации, она не зависит от idf (у схемы ltc.ltc длина документа зависела бы от idf всех его термов).
    \item Нормировка запроса на ранжирование не влияет (общий множитель), но делает оценку косинусом от 0 до 1.
\end{itemize}
Проверка на примере из лекции: документ \enquote{car insurance auto insurance}, запрос \enquote{best car insurance}, $N = 10^6$, $\mathit{df}$ = 5\,000, 50\,000, 10\,000, 1\,000 --- оценка $0{,}27 + 0{,}53 = 0{,}80$ (тест~1).

\subsection{Нечёткий и булев запрос}
\begin{itemize}
    \item \textbf{Нечёткий} --- в запросе только слова и пробелы (нет ни одного из символов \texttt{\& | ! ( )}). Найденными считаются документы хотя бы с одним словом запроса, порядок --- по оценке. Документ со всеми словами обычно выше, но короткий документ с редким словом может оказаться выше длинного со всеми словами. Разобранный запрос печатается как \texttt{any(слова)}.
    \item \textbf{Булев} --- есть хотя бы один оператор. Множество найденных документов строго то же, что в ЛР8 (пробел и \texttt{\&\&} --- И, \texttt{||} --- ИЛИ, \texttt{!} --- НЕ, скобки), порядок --- по оценке. В оценку входят только слова не под НЕ: в \texttt{rose !garden} документы упорядочены по весу \texttt{rose}. Запрос из одних отрицаний (\texttt{!lviv}) даёт всем документам оценку 0.
    \item При равной оценке выше документ с меньшим docID --- выдача однозначна.
\end{itemize}

\subsection{Изменения формата индекса}
Индексатор ЛР7 получил ключ \texttt{-t}. Формат BIDX версии 1 расширен ещё одним флагом, как ЛР4 и ЛР11:
\begin{itemize}
    \item \texttt{flags}, бит 2 (значение 4) --- \enquote{в списках есть частоты}. Флаги независимы: со сжатием (\texttt{-c -t}) --- 6, ещё и со стеммингом --- 7.
    \item Запись документа --- 16 байт вместо 12 (поле \texttt{doc\_rec} заголовка): к заголовку, URL и числу токенов добавлено поле \texttt{norm} --- длина вектора документа $\|d\|$, 4 байта, число с плавающей точкой IEEE~754 одинарной точности, порядок байтов little-endian. Старый читатель шагает по записям с размером из заголовка, поэтому лишнее поле ему не мешает.
    \item Запись списка документов --- пара $(\mathrm{docID}, \mathit{tf})$:
    \begin{itemize}
        \item без сжатия --- 8 байт: docID (4 байта) и $\mathit{tf}$ (4 байта); \texttt{first} --- номер записи, список начинается с байта $\texttt{postings\_off} + 8 \cdot \texttt{first}$;
        \item со сжатием --- коды VByte (ЛР11) парами: разность docID, затем $\mathit{tf} \ge 1$; \texttt{first} --- смещение в байтах.
    \end{itemize}
\end{itemize}
\begin{center}
\small
\begin{tabular}{|l|l|}
\hline
Флаги & Запись списка документа $d_i$ с частотой $\mathit{tf}_i$ \\
\hline
0 & \texttt{[docID u32]} \\
4 & \texttt{[docID u32][tf u32]} \\
2 & \texttt{[VByte($d_i - d_{i-1}$)]} \\
6 & \texttt{[VByte($d_i - d_{i-1}$)][VByte(tf)]} \\
\hline
\end{tabular}
\end{center}
Пример (тест~3): документы \enquote{rose rose rose garden}, \enquote{rose flower}, \enquote{garden}; список \texttt{rose} --- $(0, 3), (1, 1)$, со сжатием --- байты \texttt{80 83 81 81}; длины векторов документов --- $\sqrt{(1 + \log_{10} 3)^2 + 1} \approx 1{,}78$, $\sqrt{2}$ и 1.

Остальные читатели индекса (булев поиск ЛР8, \texttt{lookup}, \texttt{listbench} ЛР11) берут из пар только docID: метод \texttt{read\_list} читателя \texttt{index\_reader.h} пропускает частоты, а новый метод \texttt{read\_list\_tf} возвращает и их.

\subsection{Индексация}
Частота считается при вставке токена в словарь: если терм уже встречался в текущем документе (номер последнего документа терма равен текущему), увеличивается частота последней записи его списка, иначе дописывается новая запись с частотой 1. Частоты лежат в отдельном массиве рядом с массивом docID терма, поэтому без ключа \texttt{-t} память не расходуется. Длины векторов документов считаются после обработки всех документов одним проходом по всем спискам (сумма квадратов весов по каждому документу) и записываются в таблицу документов.

\subsection{Выполнение запроса}
Код --- \texttt{lab13/src/rank.h} (без STL), подключён в \texttt{build/search} ЛР8:
\begin{enumerate}
    \item Запрос разбирается тем же парсером ЛР8 (с токенизацией и, если у индекса флаг стемминга, стеммингом слов). Из дерева собираются слова не под НЕ с их частотой в запросе; считаются веса запроса.
    \item Накопители оценок: массив из $N$ чисел \texttt{double} и массив отметок \enquote{документ встретился}. Для каждого слова запроса его список распаковывается вместе с частотами, и к накопителю каждого документа списка прибавляется $w_{t,q} \cdot (1 + \log_{10} \mathit{tf}_{t,d}) / \|d\|$. Номера затронутых документов запоминаются, чтобы потом обнулить только их, а не весь массив.
    \item Нечёткий запрос: результат --- все затронутые документы. Булев: результат --- множество, вычисленное как в ЛР8, оценки берутся из накопителей.
    \item Результат сортируется слиянием по убыванию оценки, при равенстве --- по docID. Время --- $O(\sum_t \mathit{df}_t + n \log n)$, где $n$ --- число найденных документов.
\end{enumerate}
Утилита командной строки и веб-сервис ЛР8 работают как раньше, но по индексу с частотами в выдаче пятый столбец --- оценка, а веб-сервис показывает её под каждым результатом (\texttt{python3 lab8/web/server.py --index lab13/data/index-tf.bin}).

\subsection{Оценка качества}
\begin{itemize}
    \item \textbf{Запросы.} 12 запросов из слов без операторов (\texttt{data/queries.tsv}), у каждого записана информационная потребность --- критерий релевантности: \texttt{cheap hostels paris}, \texttt{best time to visit iceland}, \texttt{moscow aviation institute}, \texttt{lviv old town}, \texttt{lost passport abroad} и др.
    \item \textbf{Системы.} (1) Булев поиск ЛР8 --- пробел значит И, выдача по docID; (2) TF-IDF, тот же запрос как нечёткий; (3) TF-IDF, слова запроса через \texttt{\&\&} --- то же множество, что у (1), но упорядоченное по оценке. Сравнение (1) и (3) показывает вклад одного только порядка, (2) --- ещё и нечёткого совпадения.
    \item \textbf{Метрики} по первым 10 результатам: P@10 --- доля документов с оценкой $\ge 1$; NDCG@10 --- $\mathrm{DCG} = \sum_{i=1}^{10} (2^{g_i} - 1) / \log_2(i + 1)$, делённое на DCG идеального порядка всех размеченных документов запроса; ERR@10 --- ожидаемая обратная позиция, на которой пользователь найдёт ответ: $\sum_r \frac{1}{r} R_r \prod_{i<r} (1 - R_i)$, $R = (2^g - 1) / 4$. Если найдено меньше 10 документов, недостающие считаются нерелевантными.
    \item \textbf{Разметка.} Пул --- объединение первых 10 результатов трёх систем, 282 документа (\texttt{scripts/pool.py}). Оценки 0 --- нерелевантен, 1 --- частично, 2 --- релевантен ставились вслепую: в файлах для разметки документы перемешаны, не видно, какая система их нашла. Разметку выполняли три оценщика-субагента (языковая модель Claude), каждый --- по четырём запросам, по заголовку, началу текста и фрагментам со словами запроса, при сомнении --- по полному тексту. Оценки с обоснованием каждой --- \texttt{data/qrels.tsv}.
\end{itemize}
